\section{Payload encryption, key custody and authorized delivery}\label{sec:custody}
Every stored payload is encrypted. A bucket's encryption mode (\rr{ENC-001}, decision D-12) decides who encrypts it and who can decrypt it. In \term{server mode}, the default, the storage role encrypts each version under a fresh data key that is wrapped by a tenant key-encryption key (KEK) held only by the \code{keys} role, and decrypts for a reader only after the \code{keys} role has released that data key under a fresh signed decision for the exact operation. In \term{client mode}, which a bucket opts into, clients encrypt content before upload and decrypt it after download, and no platform role can read it. R2 MCP follows the same modes: it returns server-mode content as plaintext after the same checks and client-mode content as encrypted packages. TLS protects every connection in both modes, and payloads remain ciphertext in backing stores and backups.

For client-mode content, no platform role receives plaintext payloads, raw data keys or recipient private keys. The \code{keys} role is then an authorization and envelope-delivery service: it stores and returns opaque, client-created recipient envelopes and cannot unwrap them, rewrap one for another user, or recover plaintext. For server-mode content it holds the tenant KEKs and wraps and unwraps data keys only under signed decisions; it never sees blobs or payload plaintext. Independently authored Go code integrates qualified streaming AEAD, HPKE and signature libraries; it does not invent cryptographic primitives.

\subsection{Encryption modes}
\begin{longtable}{@{}P{0.11\textwidth}P{0.24\textwidth}P{0.30\textwidth}P{0.25\textwidth}@{}}
\caption{Encryption modes}\label{tab:modes}\\
\toprule Mode & Encrypts & Decrypts for a reader & Sharing\\\midrule\endfirsthead
\toprule Mode & Encrypts & Decrypts for a reader & Sharing\\\midrule\endhead
\code{server} (default) & The storage role, with a fresh data key per version wrapped under the tenant KEK & The storage role, after the \code{keys} role releases the data key under a fresh decision for the exact operation & Policy decides who reads, including through share grants (\rr{SHR-002}); no recipient key is needed\\
\code{client} (opt-in) & The uploading client, before upload & Only clients holding a recipient envelope and its private key & Policy allows the reader, and a key-holding client publishes the reader's envelope\\
\bottomrule
\end{longtable}
\reqn{ENC-001}{R1-B}{T-06, T-23}{Every bucket has exactly one encryption mode, \code{server} or \code{client}, set through the authorized control API when the bucket is created; \code{server} is the default. There is no plaintext mode. The mode is protected bucket metadata and is visible to policy as \code{resource.attributes.encryption} on the bucket and on every resource stored in it, including objects, document versions, units, entries and derived outputs, so that a business rule or platform ceiling can require a mode, for example for a label. Changing it is a registered operation that needs \code{manage} on the bucket and may need an approval; it applies only to uploads begun after the change commits.}
\reqn{ENC-002}{R1-B}{T-23}{Each immutable version records the mode it was published under and keeps it. A client-mode version is never read, delivered, exported or migrated through a server-held key, and a server-mode version never takes a client-created recipient envelope or client signature in place of its server binding. A route for one mode refuses a resource of the other mode with a stable error, after the existence checks of \rr{CNT-004}. A version moves to the other mode only as a new version: an authorized principal reads it through its own mode and writes it again under the ordinary write checks.}
\reqn{ENC-003}{R1-B}{REV}{Documentation, administrative listings and audit state which mode applies to each bucket and version. The guarantee that no platform role can read content is claimed only for client-mode versions; server-mode content carries non-guarantee N-09. Upgrading a deployment from a release without modes places its existing buckets in client mode, so already-published content keeps its custody.}

\subsection{Server mode: platform-held keys}
Server mode trades the client-custody guarantee for reach: any authenticated client, including a generic MCP client, can read what policy allows, and no reader needs a recipient key. Its protection is the access decision, enforced inside the authority fence at every key release, and the separation of KEKs (\code{keys} role) from blobs and plaintext (storage role).
\reqn{SSE-001}{R1-B}{T-23}{In server mode the storage role encrypts each object version, unit version, original file and derived-output version with a fresh data key under the qualified streaming AEAD of \rr{CRY-002}. The AAD binds tenant, object, logical version, crypto version, blob, suite, format and the server mode under a domain identifier distinct from the client-mode AAD of \rr{CRY-003}, so ciphertext copied to another identity or presented as client-mode content fails authentication. Data keys are never shared across tenants, objects, units or versions, and no bucket-wide or document-wide data key exists.}
\reqn{SSE-002}{R1-B}{T-18, T-23}{Each tenant has versioned key-encryption keys held only by the \code{keys} role. The catalog stores each server-mode data key only wrapped under a KEK, with the KEK version and a binding to the same identity as the AAD. No other role holds a KEK, and the \code{keys} role never holds blobs, ciphertext streams or payload plaintext. KEK rotation adds a version and rewraps data keys inside the \code{keys} role without changing ciphertext; a retired version is kept until no wrapped key refers to it. Production KEK custody is decision D-13; until it is closed, KEKs are development key files protected like other local key files, and no production payload is accepted in server mode.}
\reqn{SSE-003}{R1-B}{T-13, T-23}{The \code{keys} role generates and wraps a data key only against a signed decision for the exact write, naming its upload intent, and unwraps one only against a single-use signed decision for the exact read, export or disclosure operation and version. It verifies that decision inside the ordering fence of \rr{DAT-007} with the same bindings as recipient-envelope delivery: audience and exact check set, current identity and revisions, policy epoch, recovery incarnation, approval, the head state of a current-version read and, for an original file, every unit check of \rr{CNT-003}. A missing, stale or misbound decision releases no key. Every read obtains a fresh release, so revoked authority takes effect at the next admission without retiring any envelope. Each wrap and unwrap is audited with resource, version, operation and decision identifiers, KEK version and outcome, never key material.}
\reqn{SSE-004}{R1-B}{T-19, T-23}{An unwrapped data key exists only in the memory of the \code{keys} role for one wrap or unwrap and of the storage role for one upload or read stream. It is never persisted, logged, cached across requests or returned to a client, and it never reaches the decision, MCP, tools or maintenance roles. Key buffers have an explicit owner and lifetime and are cleared after use as far as the runtime permits (\rr{KEY-007}). Server-mode plaintext is likewise kept out of logs, traces, dumps, temporary files and support exports.}
\reqn{SSE-005}{R1-B}{T-08, T-23}{A server-mode upload is encrypted as it arrives: plaintext never reaches a staged blob, temporary file or spill. The storage role records ciphertext digest and length and obtains the durable receipt before publication under the ordinary protocol (\rr{PUB-002}, \rr{PUB-003}). A retry with the same idempotency key returns the outcome of the attempt it names and cannot change bytes already staged or published; a restarted encryption uses a fresh data key and fresh version, crypto-version and blob identities (\rr{CRY-007}). On read the storage role decrypts authenticated segments in order and aborts the connection on an authentication, digest or stream-termination failure; it never reports a shorter plaintext as complete. Range reads remain unsupported.}
\reqn{SSE-006}{R1-B}{T-18, T-23}{Backups contain wrapped data keys with the catalog, never KEKs or unwrapped keys. Restore needs the KEK versions that the restored wrapped keys refer to, recovered separately under D-13. Reconciliation verifies that each referenced KEK version is present and quarantines server-mode versions whose keys cannot be unwrapped before admission reopens (\rr{AVL-008}). Final removal of a server-mode version destroys its wrapped data key; the content is unrecoverable once no retained backup holds that wrapped key, and the platform reports that point rather than claiming it earlier.}

\subsection{Client mode: client custody}
The rest of this part applies to client-mode versions unless a requirement names server mode. Each immutable object version has a fresh data key (keyset). In R1-C each separately protected text-unit version is an ordinary immutable object with its own fresh keyset. A client encrypts the payload under that keyset and wraps the keyset to each explicitly approved recipient public key using HPKE. The server stores ciphertext and opaque recipient envelopes with authenticated bindings. There is no shared bucket data key, server master decryption key or whole-document data key that unlocks independently protected units.

\begin{longtable}{@{}P{0.27\textwidth}P{0.63\textwidth}@{}}
\toprule Boundary (client mode) & Material and responsibility\\\midrule\endhead
Owner or authorized recipient client & Plaintext, its private keys and the data keys it can unwrap; encrypts, verifies and decrypts locally\\
Optional organization custodian client & An independently controlled client with an explicit recipient envelope; can recover or authorize new recipient envelopes only for versions it can already decrypt\\
Platform storage role & Ciphertext, signed descriptors and catalog bindings; verifies publication, access and durability without decrypting\\
Platform \code{keys} role & Opaque recipient envelopes and public-key records; authorizes publication and delivery without unwrap capability\\
MCP and tool services & Authorized encrypted packages and permitted routing metadata; no plaintext conversion or private-key custody\\
Backing storage and backups & Ciphertext and protected metadata; provider encryption is additional protection, not a substitute for payload encryption\\
\bottomrule
\end{longtable}
\req{CRY-001}{R1-B}{T-10}{Client encryption is the only payload custody profile of a client-mode bucket. There, missing, unknown, server-decryptable or unencrypted profiles reject publication, and neither an API flag nor a privileged server role can request plaintext storage or server-side decryption of a client-mode version. In either mode, profile, suite, format and mode are recorded for each immutable version and checked on every operation that reads or delivers it.}

\subsection{Payload cipher suite and immutable binding}
\req{CRY-002}{R1-B}{T-10}{The client SDK, and the storage role for server-mode content, use a qualified streaming AEAD (initial candidate: Tink AES-256-GCM-HKDF with 1 MiB segments). The library handles nonce generation, key derivation, segment authentication and final-segment marking; project code \mustnot\ reproduce its chunked encryption. The module checksum and format tests are pinned.}
\req{CRY-003}{R1-B}{T-10}{Payload additional authenticated data (AAD) binds the tenant, object, crypto version, cipher suite and format under a versioned domain identifier. The same binding appears inside each recipient key package, preventing substitution across those identities. Mutable display names, classifications and envelope generation are excluded so that authorized metadata and recipient changes do not rewrite payloads. AAD is authenticated, not confidential; secrets and personal data do not belong in it. The exact field encoding is retained in the companion \code{crypto-format-profile.tex} and is qualified with the client wire contract under \rr{RCP-007}.}

\req{CRY-004}{R1-B}{T-10}{A published client-mode version binds tenant, object, logical version and crypto-version IDs; suite and format; ciphertext digest and length; uploader signing identity; and signed ciphertext descriptor. Its envelope records additionally bind recipient fingerprint and key version, envelope generation, opaque HPKE ciphertext and the origin signature. Logical versions, crypto versions, recipient-key versions and envelope generations remain distinct. A server-mode version binds the same identities, suite, format, ciphertext digest and length, the authenticated uploader and its wrapped data key and KEK version (\rr{SSE-002}) instead of a client signature and envelopes.}
\req{CRY-005}{R1-B}{T-11}{Client-mode raw keysets and recipient private keys never enter platform APIs, logs, persisted records or process memory; only authorized clients handle them. Server-mode data keys follow \rr{SSE-004} and never leave the memory of the storage and \code{keys} roles. A public key or client-supplied opaque envelope is not authority to read, create or share a resource; the server checks the exact operation independently of possession of those bytes.}

\subsection{Client-generated keys and server envelope service (client mode)}
\reqn{KEY-001}{R1-B}{T-10, T-12}{Client mode uses flat recipient wrapping: the client generates a fresh data keyset for every immutable object version and, in R1-C, every protected unit version, then creates a separate HPKE envelope for each explicitly approved recipient. Every envelope contains exactly one keyset for one immutable object or unit version and one recipient; a batch is a container of independently authorized envelopes, never a multi-unit encrypted key bundle. Data keys are not shared across tenants, objects or versions. A document-level KEK hierarchy for client-mode keys is deferred under D-09; any later proposal keeps those KEKs solely in authorized clients, preserves independent unit release, and cannot give platform roles an unwrap capability for client-mode keys. The server-mode tenant KEK of \rr{SSE-002} is a separate profile and never wraps a client-mode key.}
\req{KEY-002}{R1-B}{T-11}{The envelope service authenticates the calling client or service and requires a signed decision for the exact delivery or publication operation. It resolves recipient fingerprints and envelope identities from the canonical registry and catalog, verifies current revisions, and returns only the exact stored envelope for that recipient. It never accepts a caller-provided key-server URL, unwraps an envelope, or substitutes a broader recipient.}
\req{KEY-003}{R1-B}{T-11}{Before encryption, an upload intent binds the stable object ID and allocates fresh logical-version, crypto-version and blob IDs for that encrypted attempt. Exact retries reuse those IDs and the identical ciphertext; a restarted encryption attempt requires a new intent with fresh version, crypto-version and blob IDs. The server validates ciphertext length and digest, signed descriptor, uploader authority, recipient bindings, durable receipts and publication preconditions; it cannot verify the inner AEAD tag, plaintext format, plaintext classification or whether an opaque envelope decrypts. Those properties depend on the authorized publishing client and are checked by receiving clients where possible. Outer format validation cannot prove that a malicious uploader encrypted arbitrary submitted bytes; the confidentiality guarantee assumes a trusted compatible publishing client, not truthful labels on hostile uploads. Pending ciphertext and envelopes are not eligible for general reads.}
\req{KEY-011}{R1-B}{T-10, T-11}{The client and server operations of client mode behave as in Table~\ref{tab:broker}. Their server operation names and API types \mustnot\ imply plaintext key generation, unwrap or rewrap. The server-mode key operations of \rr{SSE-003} are separate operations, named as such, and never accept or return client-mode key material.}
\reqn{KEY-012}{R1-B}{T-11}{Ciphertext access, recipient-envelope delivery and sharing are separate compound operations with the check sets of Table~\ref{tab:keyops}. The storage and envelope services verify the complete applicable signed result, exact operation and recipient bindings; absence of a required check denies. Initial envelopes supplied with an upload pass the same publication checks as later envelopes, including recipient-revision checks at commit. Publisher share/create checks use the authenticated publisher actor and subject; target read/read-units and receive-key eligibility use the server-resolved recipient principal through \rr{RES-009}. Each signed check names its evaluated principal and current identity revision. Initial-upload eligibility is evaluated for the exact planned published state bound to the intent, not as permission to read STAGED content. An eligibility result grants no delegation or present read authority; caller-supplied recipient attributes cannot replace authoritative resolution. R1-C applies the same checks to both the document boundary and every selected unit.}
\begin{longtable}{@{}P{0.30\textwidth}P{0.47\textwidth}P{0.13\textwidth}@{}}
\caption{Compound operations for encrypted content}\label{tab:keyops}\\
\toprule Operation & Required checks and bindings & Enforcer\\\midrule\endfirsthead
\toprule Operation & Required checks and bindings & Enforcer\\\midrule\endhead
Ciphertext read & \code{read} on the exact object version; for documents, parent \code{read-units} and \code{read} on every selected unit; permitted metadata projection & Storage\\
\code{GetRecipientEnvelope} & Caller read and envelope-release checks; separate target-recipient read/read-units and \code{receive-key} eligibility under \rr{RES-009}; exact recipient fingerprint and current revisions & Envelope service\\
\code{PublishRecipientEnvelope} & Publisher \code{share} on the exact object or document and selected units, plus envelope \code{create}; separate target-recipient read/read-units and \code{receive-key} eligibility; client-origin signature and exact ciphertext binding & Envelope service\\
Remote package export & \code{export} on the object or document version and \code{export-key} on each selected unit where applicable, in addition to the applicable read and envelope-delivery checks; only an already-created envelope for that recipient & Storage and envelope service\\
\bottomrule
\end{longtable}
\begin{longtable}{@{}P{0.29\textwidth}P{0.61\textwidth}@{}}
\caption{Client and envelope-service operations}\label{tab:broker}\\
\toprule Operation & Required behavior\\\midrule\endfirsthead
\toprule Operation & Required behavior\\\midrule\endhead
\code{BeginEncryptedUpload} (server) & Authorize create or new-version; bind stable object ID and allocate intent, logical-version, crypto-version and blob IDs; return parameters without generating or receiving a data key\\
\code{EncryptObjectVersion} (client) & Generate keyset; stream-encrypt with exact AAD; create recipient envelopes for independently verified approved public keys; sign ciphertext descriptor and envelope contexts\\
\code{FinalizeEncryptedUpload} (server) & Verify signature, intent, profile, ciphertext digest/length and durable receipt; apply the full initial-envelope checks. Publish at the authorized object commit, or retain document units STAGED until the atomic document commit\\
\code{GetRecipientEnvelope} (server) & Verify the full signed delivery decision and freshness; return the existing opaque client-created envelope with its origin evidence; record delivery audit\\
\code{CreateRecipientEnvelope} (client) & Obtain a usable existing envelope and verify source; authorize sharing; verify target fingerprint independently; unwrap locally and create a new HPKE envelope for the same data key and exact version; sign its context\\
\code{PublishRecipientEnvelope} (server) & Recheck publisher authority, target-recipient eligibility and both identity and recipient revisions at commit; verify client origin and bindings; atomically publish opaque envelope, generation and audit record\\
\code{RetireEnvelope} (server) & Authorize retirement; withdraw future delivery, advance revisions and record audit; preserve required recovery references without claiming to erase client copies\\
\bottomrule
\end{longtable}

An owner or organization custodian is a client connected to the same system, not a privileged server decryption role. It can create a new recipient envelope only if it already holds the corresponding private key or a usable data key. Approval alone cannot manufacture access to an existing encrypted version. If no such client is available, the share operation remains unavailable; there is no emergency server unwrap.

\subsection{Client key custody and failure handling}
\req{KEY-004}{R1-B}{T-10, T-12}{Recipient-key rotation is performed by an authorized client: decrypt an existing envelope locally, verify the newly approved recipient key, create and sign a replacement envelope, and publish it under the expected envelope and recipient revisions. The data key and ciphertext identity remain unchanged. The platform cannot rewrap it itself, and rewrapping the same disclosed data key does not revoke old access.}
\req{KEY-005}{R1-B}{T-10, T-12}{The client uses its approved private-key store and a qualified cryptographic library. Optional hardware or an organization-operated key service belongs to the explicitly trusted client custody boundary and is selected under D-08; it is not a platform KMS unwrap dependency. No platform identity has permission to use recipient private keys. Envelope encryption is randomized; retries reuse the persisted encrypted attempt and operation identity, not convergent encryption or caller-chosen nonces.}
\req{KEY-006}{R1-B}{T-10}{Inside HPKE, a versioned key package carries the exact payload binding and one library-serialized binary keyset, with explicit bounded lengths. The receiving client rejects inconsistent lengths, trailing bytes, unknown versions, keysets above 64 KiB, extra keys or keysets, and any suite, binding or streaming-AEAD-template mismatch before use. The exact field ordering and integer encodings are retained in \code{crypto-format-profile.tex}. Servers treat the package as opaque and enforce only outer encoding and size bounds; they cannot prove that a hostile publisher put only the declared key inside HPKE.}
\req{KEY-007}{R1-B}{T-19}{Clients protect keys and plaintext against logs, traces, dumps and support exports and restrict debugging of decryption processes. Raw key buffers have an explicit owner and lifetime, avoid strings and unnecessary copies, and are cleared when the API permits. Perfect zeroization of Go or library-managed copies is not promised. Server diagnostics contain only permitted metadata and opaque-envelope identifiers; TLS protects all client and service connections in addition to payload encryption.}
\reqn{KEY-013}{R1-B}{T-11}{Process and credential boundaries prohibit all platform roles, including storage, decision, MCP, tools, maintenance and envelope delivery, from receiving client-mode payload plaintext or client-mode data keys. For server-mode content only the storage role handles plaintext and unwrapped data keys and only the \code{keys} role handles KEKs (\rr{SSE-002}, \rr{SSE-004}); the decision, MCP, tools and maintenance roles receive neither. The client verifies both a client-mode package's origin and its own recipient binding before decryption. Deployments \mustnot\ host a plaintext-processing custodian or model inside the platform and describe it as client-only encryption.}
\reqn{KEY-014}{R1-B}{T-11, T-13}{Server caches contain only ciphertext, opaque envelopes or permitted metadata, never plaintext keysets, unwrapped data keys or server-mode plaintext. Every new delivery or key release still checks current identity, recipient status, resource and envelope revisions, policy epoch and expiry. Client key retention follows the client's explicit custody policy; the platform does not claim that a client-side cache expires cryptographically when a grant or server lease expires.}
\reqn{KEY-015}{R1-B}{T-11, T-13}{A missing recipient envelope, unavailable client private key, failed signature or AEAD validation, stale authority, or unavailable authorization dependency prevents the corresponding delivery or decryption. A client-mode version never falls back to a server-held key (\rr{ENC-002}), another recipient, an older weaker profile or plaintext transfer, and a server-mode read whose key release fails returns no content. Client-side key-store failure does not cause private keys to be uploaded to the platform.}

\subsection{Client decryption, streaming and copy}
\req{CRY-006}{R1-B}{T-10}{The client releases plaintext only after verifying its signed origin and exact ciphertext digest and length. The baseline therefore receives the complete ciphertext into bounded, access-restricted encrypted-content storage before decrypting it as authenticated segments; it rejects oversized objects, insufficient temporary capacity, incomplete transfers and digest failures without releasing plaintext. Temporary ciphertext is cleaned up under the client custody policy. An alternative streaming profile needs equivalent authenticated origin binding and separate qualification. Overall success additionally requires valid AEAD stream termination. Wrong AAD or key, corruption, segment reordering, duplication, truncation and forged final segments fail. Memory and concurrent work remain bounded. Qualification includes temporary-storage demand and time to first plaintext; server leases govern delivery, not subsequent use of retained bytes or keys.}
\req{CRY-007}{R1-B}{T-10}{Retrying an upload resends identical ciphertext and envelopes for the same encrypted attempt. Restarting encryption requires a new intent and fresh logical-version, crypto-version and blob identities, together with a fresh keyset; it cannot overwrite an earlier attempt or change the bytes associated with an existing idempotency identity. Server-mode retries follow \rr{SSE-005}. The baseline reads complete streams in both modes. Ciphertext range transfer and client-side authenticated range assembly require a separately qualified profile, and the server implements no plaintext ranges by decrypting.}
\req{CRY-008}{R1-B}{T-10}{AAD binds the logical object and crypto version, so ordinary byte copy to a new object is not a valid encrypted object. An authorized client decrypts locally and re-encrypts under a fresh destination key and binding, then publishes a new version; otherwise copy rejects. A server-mode copy is an authorized read and a new version written through the storage role under a fresh data key. Storage repair may copy identical ciphertext for the same immutable identity. A future alias feature needs explicit same-ciphertext semantics and cannot weaken identity binding.}

\subsection{Rotation, compromise and revocation}
\req{KEY-008}{R1-B}{T-12}{The platform and authorized clients respond to each event as in Table~\ref{tab:rotation}.}
\begin{longtable}{@{}P{0.24\textwidth}P{0.66\textwidth}@{}}
\caption{Key events and required responses}\label{tab:rotation}\\
\toprule Event & Required response\\\midrule\endfirsthead
\toprule Event & Required response\\\midrule\endhead
Recipient-key rotation & Owner or custodian client creates envelopes for independently verified replacement keys; publish with revision checks; retain recovery access according to explicit custody policy\\
Content-signing-key rotation & Retire signing authority while preserving trusted historical verification bindings; current access remains separately authorized (\rr{RCP-011})\\
Content-signing-key compromise & Stop new publications; quarantine affected origin evidence until reviewed under \rr{RCP-011}; never trust a claimed signing time alone\\
Data-key compromise & A fresh data key and ciphertext version, generated by an authorized client in client mode or by the storage role in server mode; new envelopes or rewrapping over the old key do not repair disclosure; retained old ciphertext remains exposed\\
KEK rotation (server mode) & The \code{keys} role adds a KEK version and rewraps data keys without changing ciphertext; a retired version stays until no wrapped key refers to it (\rr{SSE-002})\\
KEK, \code{keys}-role or storage-role compromise (server mode) & Rotate KEKs and treat every server-mode version whose wrapped key, data key or plaintext was reachable as disclosed; re-encrypt under fresh data keys as new versions where needed; client-mode content is unaffected\\
Recipient private-key compromise & Revoke future deliveries and new sharing to that key; assess every historical envelope recoverable with it; use fresh data and recipient keys for replacement versions where needed\\
Grant or user revocation & Block new ciphertext fetches, envelope deliveries and server-mode key releases at the effective revision; plaintext, keys and ciphertext already obtained cannot be recalled\\
Deletion or crypto-erasure & Track ciphertext, envelopes, wrapped keys, replicas, backups and retention; final removal of a server-mode version destroys its wrapped data key (\rr{SSE-006}); server deletion cannot prove destruction of keys or copies held by clients\\
Lost client keys (client mode) & Use only an explicitly provisioned recovery or custodian recipient; if every usable private key and recovery path is lost, the content is unrecoverable\\
\bottomrule
\end{longtable}
\req{KEY-009}{R1-B}{T-11}{Envelope publication, delivery, rotation and retirement record resource and version, actor and subject, decision and operation IDs, policy epoch, recipient fingerprint, envelope generation, origin signer and outcome, without plaintext or keys. Delivery audit proves release of encrypted material, not that a recipient decrypted or read it.}
\req{KEY-010}{R1-B}{T-13}{Revocation blocks new server-mediated fetches, envelope deliveries and server-mode key releases, including retries and cached responses after authority expires. It does not erase client private keys, exported data keys, retained ciphertext, plaintext or derived answers. Expiry in a signed delivery receipt is not a self-destructing key. Static-recipient HPKE does not promise forward secrecy after compromise of the recipient private key. Re-encrypting future versions limits future exposure but does not repair already-copied material.}

\subsection{Atomic lifecycle, format migration and recovery}
\reqn{KEY-016}{R1-B}{T-10, T-12}{The catalog binds one exact encryption mode, profile, payload crypto version, suite and format per immutable version. Publishing a replacement recipient envelope verifies the authorized client signature and expected recipient and envelope revisions, then atomically updates the envelope generation and audit record. Replacement or retirement in one document context does not replace or retire envelopes required by another retained document version. Client-mode re-encryption creates fresh keys and an immutable ciphertext version in the client, followed by the ordinary publication checks. A partial rotation cannot publish an envelope bound to another ciphertext or recipient.}
\reqn{CRY-009}{R1-B}{T-10, T-12}{Migration between qualified client-encryption formats is an authorized client read-and-new-version operation. The client verifies the source, decrypts locally, encrypts with fresh state and publishes signed descriptors and new recipient envelopes. The original remains current until the new version commits. R1-C document migration publishes a new manifest only when every referenced unit version is ready. There is no plaintext profile to migrate. A server-mode version moves to a new format as a new version written by the storage role under the ordinary write checks, and moving a version between modes follows \rr{ENC-002}; there is no in-place server-side conversion.}
\reqn{CRY-010}{R1-B}{T-11, T-13}{Migration preserves tenant and logical resource identity, labels, current access controls and historical restrictions. The migrating client receives only explicit source-read, destination-publication and share authority. The catalog records the profile of every retained version and migration coverage; tightening the minimum accepted profile rejects new delivery of nonconforming versions. Clients cannot request a weaker format as a fallback, and a retained older representation is not declared erased or migrated.}
\reqn{KEY-017}{R1-B}{T-12, T-18}{Recovery restores a mutually consistent catalog, signed descriptors, ciphertext, opaque recipient envelopes, wrapped server-mode data keys and historical content-signing trust records; KEKs are recovered separately under \rr{SSE-006}. It verifies identities, digests, signatures and references and quarantines missing or mismatched records. Admission reopens only under the recovered-authority rules of \rr{AVL-008}; restoring old records cannot silently reactivate revoked authority. A separately controlled authorized client demonstrates client-mode decryption using its own private key or an explicitly provisioned organization-custodian envelope; the restore service never receives that key or plaintext. Restore qualification covers retained versions across recipient and content-signing-key rotation, compromised-signer handling, KEK rotation, and continued denial of unauthorized server deliveries. Client recovery custody and private-key backups remain separate; platform backups alone cannot recover lost client private keys.}

\subsection{Recipient enrollment, HPKE and authenticated origin (client mode)}
Online policy checks govern ciphertext and envelope delivery; recipient-bound HPKE ensures that those server-held envelopes are usable only by holders of the corresponding private keys. Neither mechanism controls an authorized recipient's later copying or independent sharing of plaintext. The owner must explicitly choose any organization recovery recipient; there is no implicit platform escrow.

\req{RCP-001}{R1-B}{T-03, T-11}{Recipient public-key encryption, client content signing, request authentication and server decision signing are separate key purposes. Signing or OAuth proof keys are not interchangeable with X25519 HPKE encryption keys. Each registry record binds tenant, principal or workload, fingerprint, purpose, suite, validity, status and revocation revision. A public key identifies an encryption recipient but does not itself grant access.}
\req{RCP-002}{R1-B}{T-03, T-11}{Enrollment and rotation use an authenticated authorized operation and validate public-key format and suite. Activation requires trusted provisioning or proof of the private key through a server-generated, short-lived, one-use challenge bound to identity and fingerprint, completed over the authenticated session. Before encrypting or sharing, the publishing client additionally verifies the recipient fingerprint through owner-approved provisioning or an independently trusted organization-signed directory and pins that binding. Possession proof and a platform-served directory alone do not establish the intended recipient: a compromised server could substitute a key it owns. Recipient rotation therefore requires explicit verification and cannot silently replace a pinned key.}
\req{RCP-003}{R1-B}{T-03}{Issuer verification keys come only from configured trusted discovery or JWKS endpoints with bounded refresh, redirects and egress and with issuer, audience, purpose and algorithm checks. Caller-supplied \code{jku}, \code{x5u} and arbitrary key-server URLs are never dereferenced. Issuer-key discovery is not a recipient directory or authority to replace a client's pinned recipient or content-signing key.}
\req{RCP-004}{R1-B}{T-10, T-11}{The client creates recipient envelopes using a qualified Tink hybrid-encryption implementation with HPKE suite \code{DHKEM_X25519_HKDF_SHA256_}\allowbreak\code{HKDF_SHA256_AES_256_GCM}. HPKE encrypts the small data-key package, not the bulk object. It does not by itself authenticate the author, so client-origin signatures and trusted signer bindings are required as in \rr{RCP-008}.}
\req{RCP-005}{R1-B}{T-11}{The published recipient envelope binds exact tenant, object and logical and crypto versions, ciphertext digest and length, profile, recipient fingerprint and key version, envelope generation and authorized sharing client. For a document unit, its identity additionally includes the document, document version and opaque manifest digest; distinct contexts coexist even when they reference the same immutable unit and recipient. It contains exactly that one unit's keyset; independently authorized envelopes for multiple units may be returned together but are never combined inside one HPKE ciphertext. The envelope service returns stored envelopes only after all applicable checks of Table~\ref{tab:keyops}; it cannot manufacture a missing recipient envelope or expand section-only rights to other units.}
\req{RCP-006}{R1-B}{T-10}{The client serializes a versioned Protobuf \code{RecipientEnvelopeContextV1} once under the explicit application wire profile of \rr{RCP-007}, with fixed UUID and byte representations, timestamp units and exactly one object-or-unit version and recipient binding. It signs those exact context bytes with the HPKE ciphertext; receivers verify received bytes rather than reserializing them. Protobuf serialization is not assumed canonical~\cite{protobufCanonical}. Strict wire validation rejects unknown or duplicate fields, aggregate bindings, unsupported values and oversized fields before a permissive decoder can discard ambiguity. HPKE context information is a domain-separated SHA-256 of the context bytes. A batch is a bounded, uniquely ordered list of separate envelopes, each independently validated and authorized. Stable encrypted context contains no expiring decision: delivery carries a separate signed receipt bound to envelope digest, operation, decision, audience, current authority and recovery revisions, and expiry.}
\req{RCP-007}{R1-B}{T-10}{Before the encrypted-package profile is enabled, its versioned interoperability contract freezes the message schemas, parsing rules and limits, signature and digest algorithms, domain separation, fingerprint encoding, Tink framing and key encoding, and exact positive and negative wire examples. Actual client/server interoperability tests qualify that contract and the pinned library versions; naming Protobuf, Tink or HPKE alone is not qualification. Algorithm or framing changes require a new reviewed profile rather than an implicit fallback. The native encrypted-package client is required in R1-B. An R2 MCP client needs a compatible local bridge or SDK to decrypt; links point only to authenticated delivery endpoints.}
\req{RCP-008}{R1-B}{T-10, T-11}{The publishing client signs the ciphertext descriptor and each recipient-envelope context and ciphertext with a separately registered content-signing key. Receivers verify origin against independently trusted signer bindings, applying \rr{RCP-011}, and verify the ciphertext digest, exact AAD and recipient binding before releasing plaintext under \rr{CRY-006}. A document unit may use a signed descriptor bound to the opaque full-manifest digest without exposing denied entries. The manifest digest covers its explicitly encoded core before descriptors or signatures are added, avoiding self-reference. A server decision signature attests authorization and delivery metadata, not client content origin. Client software distribution and updates preserve this boundary; compromised client software can expose plaintext.}
\req{RCP-009}{R2}{T-03}{Sender-constrained OAuth such as DPoP is an optional negotiated transport profile using a maintained implementation and enforcing signature, method, URI, token-key binding, time, nonce and replay rules. It neither grants resource permission nor automatically binds the request body or separate HPKE recipient key. It is not required for baseline client encryption, and base MCP support is not assumed.}
\req{RCP-010}{R1-B}{REV}{ABE, proxy re-encryption and client KEK hierarchies are deferred. A later proposal needs its own threat model, maintained-library qualification and measured need. It cannot introduce platform plaintext or private-key custody for client-mode content, and it is not a parallel default implementation. The server-mode tenant KEK of \rr{SSE-002} is not such a hierarchy and never wraps client-mode keys.}
\reqn{RCP-011}{R1-B}{T-10, T-12, T-18}{Content-signing keys have immutable identity and purpose bindings and distinct active, retired and compromised states. Routine rotation or principal removal prevents new signing authority but preserves independently trusted historical verification records for retained content; it grants no current read access. Compromise blocks new publications signed by that key and causes affected historical signatures to fail closed unless a reviewed, object-specific disposition establishes trustworthy prior publication or fresh endorsement by an authorized client. An uploader-supplied timestamp is not that evidence. Clients apply authenticated trust updates with a declared freshness policy, and offline verification is not claimed to establish current revocation status. Recovery preserves these records and never converts a compromised key back to active.}
\begin{longtable}{@{}P{0.23\textwidth}P{0.43\textwidth}P{0.20\textwidth}@{}}
\caption{Cryptographic choices}\label{tab:abe}\\
\toprule Approach & Fit and limit & Decision\\\midrule\endfirsthead
\toprule Approach & Fit and limit & Decision\\\midrule\endhead
Server AEAD under a tenant KEK & Independent version keys wrapped by a KEK in the \code{keys} role; any authorized client reads without key custody; the platform can read the content, and the access decision at every key release is the protection & Default (server mode, D-12)\\
Client AEAD plus recipient HPKE & Independent version keys, opaque server storage, explicit client-managed sharing; retained keys remain usable offline & Opt-in (client mode)\\
TLS on every connection & Protects authentication, transport metadata and server-mode plaintext in transit; does not replace payload encryption & Required\\
Client or organization custodian & Recovery and new sharing only for versions with an explicit usable custodian envelope & Optional, owner-approved\\
Ciphertext-policy ABE & Offline attribute-key infrastructure does not directly implement arbitrary dynamic policies or revoke retained keys & Deferred\\
Proxy re-encryption & Could transform recipient envelopes without server plaintext, but adds another cryptographic lifecycle and trust model & Deferred\\
Client document KEK & May reduce envelope work but broadens compromise impact; may never become a platform-held master key & Deferred under D-09\\
\bottomrule
\end{longtable}
